%% 作者简历及在学研究成果
\chapter*{作者简历及在学研究成果}\addcontentsline{toc}{chapter}{作者简历及在学研究成果}

\section*{一、主要教育经历/工作经历（从大学起，到硕士入学止）}


% ...
\noindent\begin{tabularx}{\textwidth}{|>{\centering\arraybackslash}X|>{\centering\arraybackslash}X|>{\centering\arraybackslash}X|}
  \hline
  起止年月 & 学习或工作单位 & 备注 \\\hline 
  2019年9月至2023年6月  & 在北京科技大学计算机科学与技术专业攻读学士学位  & X  \\\hline

  & & \\ \hline
\end{tabularx}

\section*{二、在学期间从事的科研工作}

% （1）高性能低功耗深度神经网络处理器芯片设计（USTB06500093），主要参与人员，2017.11-2022.11。

% （2）... 

\section*{三、在学期间所获的科研奖励}

（1）硕士研究生学业奖学金

\section*{四、在学期间发表的论文}

% [1] 	\textbf{Liu M. K.}, Liu C. X., Zhang S. T., et al. Research on Industry Development Path Planning of Resource-Rich Regions in China from the Perspective of “Resources, Assets, Capital” [J]. Sustainability, 2021, 13(7): 3988-3988.


% {\color{blue} 盲审说明（正式写作请删除此说明）：}

% 盲审论文需要隐藏掉所有会影响盲审结果的论文作者及其导师的信息，以便论文评
% 阅人能够公正的进行评阅。

% 当学校要求提供盲审论文时，请按如下方法制作。

% 1) 对于论文中下列学生和导师信息。请将学生姓名、学生学号、导师姓名，依次全
% 部替换为[本论文作者] 、[论文作者学号]、[本论文导师]。论文中上述信息均需要替换，
% 包括作者研究成果等部分的有关信息。

% 2) 在研究成果中，论文作者发表的文章列表中应隐去所有作者的名字，只标明论文
% 期刊名，级别，发表年份。

% 3）盲审论文，请不要填写致谢，致谢页除标题、页眉、页码外请保持空白。

% 4）其他会影响盲审结果的信息，请采用类似方式处理。

% 5）提交的盲审论文应为正式论文，除了上述替换后的信息外，应为可以评阅的正式
% 论文。

% 6）论文封面，请填写专业名称、论文题目等，将学号和姓名项目保持空白不填。制
% 作盲审论文时，论文书脊、封二、题名页请删除。

% 替换前后将文档分别保存，以便盲审论文与其他论文分开管理。

% {\color{blue} 盲审样例（正式写作请删除此样例）：}

% 五、在学期间发表的论文：

% [2] \textbf{第二作者（导师一作）}.中国矿业.中文核心期刊.2020.


\pagestyle{empty}